\documentclass[12pt]{article}

% ---------- Compatibility & Fonts ----------
\usepackage[utf8]{inputenc}
\usepackage[T1]{fontenc}
\usepackage{lmodern}

% ---------- Layout & Math ----------
\usepackage[a4paper,margin=1in]{geometry}
\usepackage{amsmath,amssymb,mathtools}
\usepackage{setspace,enumitem,microtype}
\setstretch{1.2}
\setlength{\parindent}{0pt}
\setlength{\parskip}{6pt}

\begin{document}

\begin{center}
{\Large \textbf{Solutions -- Midterm Exam 2}} \\[0.9em]
\textbf{ECON 320 -- Introduction to Mathematical Economics} \\[0.7em]

November 05, 2025
 \\[0.8em]
Instructor: Muhebullah Karimzada
\end{center}

\hrule

\hrule
\vspace{8pt}

\section*{Problem 1  (12 points)}

A firm's profit: \(\pi(Q)=120Q-6Q^2-180\).

\textbf{(a) First-order condition and critical point (4 pts).}
\[
\pi'(Q)=120-12Q \quad\Rightarrow\quad \pi'(Q)=0 \Rightarrow 120-12Q=0 \Rightarrow Q^\star=10.
\]

\textbf{(b) Second derivative and classification (3 pts).}
\[
\pi''(Q)=-12<0 \quad\Rightarrow\quad \text{strictly concave} \Rightarrow Q^\star=10 \text{ is a (global) maximum.}
\]

\textbf{(c) Maximum profit (3 pts).}
\[
\pi^\star=\pi(10)=120(10)-6(10)^2-180=1200-600-180=420.
\]

\textbf{(d) Interpretation (2 pts).}
\(Q^\star\) is the profit-maximizing output. Concavity (\(\pi''<0\)) means marginal profit is decreasing; the unique stationary point is the global maximizer.

\hrule
\vspace{8pt}

\section*{Problem 2 (12 points)}

Population \(P(t)=200 e^{0.03t}\), technology \(A(t)=e^{0.02t}\), output \(Y(t)=A(t)\,P(t)^{0.5}\).

\textbf{(a) Instantaneous growth rate of \(Y(t)\) via log-differentiation (4 pts).}

\[
\ln Y(t)=\ln A(t)+\tfrac{1}{2}\ln P(t)=0.02t+\tfrac{1}{2}\big(\ln 200+0.03t\big)
= \tfrac{1}{2}\ln 200 + (0.02+0.015)t.
\]
Therefore
\[
\frac{d}{dt}\ln Y(t)=0.035 \quad\Rightarrow\quad \frac{1}{Y(t)}\frac{dY}{dt}=0.035.
\]
So \(Y\) grows at \(3.5\%\) per unit time; equivalently \(\dfrac{dY}{dt}=0.035\,Y\).

\textbf{(b) Compute \(Y(10)\) (3 pts).}
\[
Y(t)=\sqrt{200}\,e^{0.035 t}\quad\Rightarrow\quad Y(10)=\sqrt{200}\,e^{0.35}.
\]
Using \(\sqrt{200}\approx 14.1421\) and \(e^{0.35}\approx 1.4191\),
\[
Y(10)\approx 14.1421\times 1.4191 \approx 20.07.
\]

\textbf{(c) Higher population growth: \(0.03 \to 0.04\) (3 pts).}
\[
\frac{1}{Y}\frac{dY}{dt}=\frac{d}{dt}\ln Y = 0.02+\tfrac{1}{2}(0.04)=0.02+0.02=0.04.
\]
Thus the new growth rate is \(4\%\).

\textbf{(d) Interpretation (2 pts).}
Output growth equals TFP growth (2\%) plus labor/population growth weighted by its elasticity (0.5). A 1 percentage-point increase in population growth raises the output growth rate by 0.5 percentage points, showing the relative roles of technology and population.

\hrule
\vspace{8pt}

\section*{Problem 3 (13 points)}

\(f(x,y)=12x+8y-x^2-y^2-xy\).

\textbf{(a) First-order conditions and stationary point(s) (4 pts).}
\[
f_x=12-2x-y,\qquad f_y=8-2y-x.
\]
Solve \(f_x=f_y=0\):
\[
\begin{cases}
12-2x-y=0\\
8-2y-x=0
\end{cases}
\Rightarrow
\begin{cases}
y=12-2x\\
8-2(12-2x)-x=0
\end{cases}
\Rightarrow -16+3x=0 \Rightarrow x^\star=\dfrac{16}{3},\quad
y^\star=12-2\cdot\dfrac{16}{3}=\dfrac{4}{3}.
\]

\textbf{(b) Hessian and determinant (4 pts).}
\[
H=\begin{bmatrix}
f_{xx} & f_{xy}\\
f_{yx} & f_{yy}
\end{bmatrix}
=
\begin{bmatrix}
-2 & -1\\
-1 & -2
\end{bmatrix},\qquad \det(H)=(-2)(-2)-(-1)^2=4-1=3>0.
\]

\textbf{(c) Classification (3 pts).}
Since \(f_{xx}=-2<0\) and \(\det(H)>0\), \(H\) is negative definite. Therefore \((x^\star,y^\star)\) is a strict local maximum. Because \(H\) is constant and negative definite everywhere, it is also the global maximizer.

\textbf{(d) Economic interpretation (2 pts).}
Negative diagonals indicate diminishing returns to each input; \(f_{xy}=-1<0\) shows inputs are gross substitutes (raising one reduces the marginal product of the other). Global concavity implies a unique optimum.

\hrule
\vspace{8pt}

\section*{Problem 4 (13 points)}

Maximize \(U(x,y)=x^{0.5}y^{0.5}\) subject to \(2x+y=100\).

\textbf{(a) Lagrangian and FOCs (4 pts).}
\[
\mathcal{L}(x,y,\lambda)=x^{1/2}y^{1/2}-\lambda(2x+y-100).
\]
\[
\frac{\partial \mathcal{L}}{\partial x}=\tfrac{1}{2}x^{-1/2}y^{1/2}-2\lambda=0,\quad
\frac{\partial \mathcal{L}}{\partial y}=\tfrac{1}{2}x^{1/2}y^{-1/2}-\lambda=0,\quad
2x+y-100=0.
\]

%\textbf{(b) Solve for \(x^\*,y^\*,\lambda^\*\) (5 pts).}
From the first two FOCs, equate:
\[
\tfrac{1}{2}x^{-1/2}y^{1/2}=2\lambda,\qquad x^{1/2}y^{-1/2}=\lambda
\ \Rightarrow\ 
\frac{\tfrac{1}{2}x^{-1/2}y^{1/2}}{x^{1/2}y^{-1/2}}=1
\ \Rightarrow\ 
\tfrac{1}{2}\cdot\frac{y}{x}=1 \Rightarrow \frac{y}{x}=2.
\]
Thus \(y=2x\). Budget: \(2x+y=4x=100 \Rightarrow x^\star=25,\ y^\star=50\).
Multiplier from \(\lambda=\tfrac{1}{2}x^{1/2}y^{-1/2}\):
\[
\lambda^\star=\frac{1}{2}\sqrt{\frac{x^\star}{y^\star}}
=\frac{1}{2}\sqrt{\frac{25}{50}}=\frac{1}{2\sqrt{2}}\approx 0.3536.
\]

\textbf{(c) Bordered Hessian check (2 pts).}
Second derivatives:
\[
U_{xx}=-\tfrac{1}{4}x^{-3/2}y^{1/2},\quad
U_{yy}=-\tfrac{1}{4}x^{1/2}y^{-3/2},\quad
U_{xy}=\tfrac{1}{4}x^{-1/2}y^{-1/2}.
\]
Constraint \(g(x,y)=2x+y-100\) gives \(g_x=2,\ g_y=1\). The bordered Hessian
\[
H_B=\begin{bmatrix}
0 & 2 & 1\\
2 & U_{xx} & U_{xy}\\
1 & U_{xy} & U_{yy}
\end{bmatrix}_{(x^\star,y^\star)}
\]
has \(\det(H_B)>0\) at \((25,50)\), confirming a local maximum. Alternatively, note \(U\) is strictly concave on \(\mathbb{R}^2_{++}\) and the constraint is linear, implying a unique global maximum.

\textbf{(d) Interpreting \(\lambda^\star\) and envelope estimate (2 pts).}
With budget \(2x+y=I\) and \(I=100\), the envelope theorem gives \(\partial U^\star/\partial I=\lambda^\star\), the marginal utility of income. For \(\Delta I=5\),
\[
\Delta U \approx \lambda^\star \Delta I = \frac{1}{2\sqrt{2}}\times 5 \approx 1.7678.
\]

\end{document}
